% =====================================================================
%  Optimize Against EA — Handbook / User Manual
%  GameLab 2 (2025) — Julius-Maximilians-Universität Würzburg
%  Authors: Philipp Sehne, David Ahrens
%
%  Build:  pdflatex handbook.tex   (run twice for TOC/links)
%  Screenshots live in ./screenshot/.
% =====================================================================
\documentclass[11pt,a4paper]{report}

\usepackage[utf8]{inputenc}
\usepackage[T1]{fontenc}
\usepackage[english]{babel}
\usepackage[margin=2.5cm]{geometry}
\usepackage{graphicx}
\usepackage{booktabs}
\usepackage{tabularx}
\usepackage{enumitem}
\usepackage{xcolor}
\usepackage{amsmath}
\usepackage[colorlinks=true, linkcolor=blue!50!black, urlcolor=blue!50!black]{hyperref}

\graphicspath{{screenshot/}}

% --- small helpers ----------------------------------------------------
\newcommand{\todo}[1]{\textcolor{red!70!black}{\textbf{[TODO: #1]}}}
\newcommand{\route}[1]{\texttt{#1}}          % app route, e.g. \route{/PeakFinder}
\newcommand{\uielem}[1]{\textsf{\textbf{#1}}} % UI element, e.g. \uielem{Start EA}
\newcommand{\key}[1]{\fbox{\small\texttt{#1}}} % keyboard key

% Screenshot figure: \screenshot{file}{width-fraction}{caption}
\newcommand{\screenshot}[3]{%
  \begin{figure}[htbp]
    \centering
    \includegraphics[width=#2\textwidth]{#1}
    \caption{#3}
  \end{figure}}

\title{\Huge\textbf{Optimize Against EA}\\[1ex]
       \LARGE Handbook \& User Manual\\[2ex]
       \large GameLab~2 --- Julius-Maximilians-Universität Würzburg}
\author{Philipp Sehne \\ \small\texttt{Philipp.Sehne@stud-mail.uni-wuerzburg.de}
   \and David Ahrens  \\ \small\texttt{David.Ahrens@stud-mail.uni-wuerzburg.de}}
\date{\today}

\begin{document}

\maketitle
\tableofcontents

% =====================================================================
\chapter{Introduction}
% =====================================================================

\section{What is Optimize Against EA?}

\emph{Optimize Against EA} is an educational web application that teaches
optimization concepts --- in particular \textbf{Evolutionary Algorithms (EAs)}
--- through a collection of interactive mini-games. Instead of reading about
selection, crossover and mutation in the abstract, players \emph{experience}
them: they compete against an EA, watch populations evolve step by step, and
experiment with the algorithm's parameters to see how behaviour changes.

The application contains three main mini-games, each highlighting a different
facet of evolutionary optimization:

\begin{description}[style=nextline]
  \item[Peak Finder (\route{/PeakFinder})]
    Probe a hidden function surface to find its global minimum --- solo, or
    in competition with an EA whose every generation can be replayed and
    dissected step by step.
  \item[Maze Explorer (\route{/MazeExplorer})]
    Solve a fogged maze the way an EA would --- by hand-writing a fixed-length
    string of moves --- then hand the maze to a real EA, watch its population
    evolve, and discover how the choice of \emph{fitness function} (Manhattan
    distance, geodesic distance, length penalty, novelty search) completely
    changes the algorithm's behaviour.
  \item[Shooter Game (\route{/ShooterGame})]
    Fight waves of enemy agents whose behaviour is encoded as DNA and evolved
    by a genetic algorithm between rounds --- the enemies literally adapt to
    \emph{your} play style.
\end{description}

\section{Who is this manual for?}

This handbook is aimed at \textbf{players and instructors}: it explains how to
start the application, how to play each game, and what each control and
setting means. No programming knowledge is required.

Developers looking for the programmatic interfaces (engine functions, type
contracts, worker protocols) should consult the separate
\textbf{API documentation} instead.

\section{Educational goals}

Playing through the games, a player encounters the following concepts:

\begin{itemize}
  \item search spaces, local vs.\ global optima, deceptive landscapes;
  \item the evolutionary loop: evaluation $\to$ selection $\to$ crossover
        $\to$ mutation $\to$ new generation;
  \item elitism and premature convergence;
  \item fitness shaping --- how the \emph{same} EA succeeds or fails
        depending on how fitness is defined;
  \item novelty search --- why \emph{not} chasing the goal can outperform
        chasing it;
  \item continuous genomes (points on a surface), discrete genomes
        (move sequences), and behavioural genomes (agent DNA).
\end{itemize}

% =====================================================================
\chapter{Getting Started}
% =====================================================================

\section{Playing online (recommended)}

The easiest way to use the application is the live deployment:

\begin{center}
  \large\url{https://OptimizeAgainstEA.web.app}
\end{center}

No installation is required. The app runs entirely in the browser
(client-side); a modern browser is the only prerequisite.

\section{System requirements}

\begin{itemize}
  \item \textbf{Browser:} a current version of Chrome, Firefox, Edge or Safari.
  \item \textbf{Input:} mouse + keyboard \emph{or} touch screen. No gamepad
        is required. All games are playable on desktop and on mobile/touch
        devices (the shooter provides a virtual joystick and aim zone on
        touch screens); a desktop browser offers the most comfortable
        experience on very small screens.
  \item \textbf{Network:} only needed to load the page (and for the
        community Raidboss feature, which uses Firebase).
\end{itemize}

\section{Running locally from source}

For local use (e.g.\ offline demos), the app can be built from source.
Requirements: \href{https://nodejs.org/}{Node.js} (developed with v24.x;
v20+ should work) and \texttt{npm}.

\begin{enumerate}
  \item Open a terminal in \texttt{OptimizeAgainstEA/OptimizeAgainstEAApp}.
  \item Install dependencies: \texttt{npm install}
  \item Start the development server: \texttt{npm run dev}\\
        Vite prints a local URL (default \url{http://localhost:5173});
        open it in your browser.
\end{enumerate}

To create a production build instead, run \texttt{npm run build}; the static
output in \texttt{dist/} can be served by any web server
(\texttt{npm run preview} serves it locally).

\section{First steps in the app}

\begin{enumerate}
  \item The \textbf{home page} (\route{/}) welcomes you and leads to the
        \textbf{Dashboard} (\route{/Dashboard}), the central game-selection
        screen.
  \item Pick a game. Each game page has its own mode selector and settings.
  \item Look for the \uielem{?}~help buttons and the \uielem{Hints} toggle
        (light-bulb icon) --- see Section~\ref{sec:help-system}.
\end{enumerate}

\screenshot{Dashboard.png}{1.0}{The Dashboard --- entry point to all
mini-games.}

\section{The built-in help and hint systems}
\label{sec:help-system}

The application ships with two layers of in-app assistance:

\begin{description}[style=nextline]
  \item[Help dialogs (\uielem{?} buttons)]
    Open a modal with two tabs: \uielem{Gameplay} (how to play) and
    \uielem{Technical} (what the algorithm is doing). Available on the
    shooter modes and other game screens.
  \item[Contextual hints (\uielem{light-bulb} toggle)]
    When enabled, short coach-marks and toasts appear at the right moment
    during play --- e.g.\ after your first probe in Peak Finder, a hint
    explains that the EA has just advanced a generation. Hints can be
    switched off at any time and re-armed with the \uielem{Reset} button.
    Your preference is remembered between visits.
\end{description}

% =====================================================================
\chapter{Peak Finder}
% =====================================================================
\label{ch:peakfinder}

\section{Concept}

Peak Finder (\route{/PeakFinder}) is a search-space exploration game in the
spirit of \emph{Battleships}. A function surface is hidden from you: every
point on the map has a value between $0$ (the global minimum --- the point
you are looking for) and $1$ (far away from any minimum). The surface also
contains \emph{local} minima --- traps that look promising but are not the
true goal.

You explore by placing \textbf{probes}. Each probe reveals the surface in a
small circle around it (as heat map), giving you partial
information to plan your next probe. You win by placing a probe inside the
\textbf{win radius} around the global minimum.

\screenshot{PeakFinderPlayMode.png}{0.62}{Peak Finder: probing a hidden
surface. The heatmap is only revealed near your probes --- blue basins are
promising, the ringed probes mark the best readings so far.}

\section{Game modes}

Peak Finder has three modes, selectable at the top of the page:

\subsection{Play}

The classic solo mode. Load a map (a random one, a saved one, or one shared
via code --- see Section~\ref{sec:mapcodes}), then click/tap anywhere to
probe. The interface shows your probe history and your best probe so far.
Find the global minimum in as few probes as you can.

\begin{itemize}
  \item \textbf{Probe:} click / tap a location on the map.
  \item \textbf{Visualization:} \uielem{Heatmap} display of the revealed areas.
  \item \textbf{Win:} a probe inside the win radius triggers the win
        overlay, which shows your statistics.
\end{itemize}

\subsection{Create}

Build your own maps. First pick a \textbf{map size} (Small, Medium, Large or
Huge) --- it determines how many minima the map holds (from 4--6 on Small up
to 20--26 on Huge), the win radius, and the minimum spacing between minima,
which is enforced while you place them. Scatter your local minima, then
choose which one is the \emph{global} minimum. Dashed rings preview each
minimum's zone. Finished maps can be saved locally and exported as a
\textbf{map code} to share with others.

\subsection{Vs EA}

The centerpiece mode: you and an Evolutionary Algorithm search the
\emph{same} hidden map --- and this time it is a race.

\begin{itemize}
  \item You probe as usual. For every probe you place, the EA advances
        exactly one \textbf{generation} --- your move budget and its
        evolution budget are coupled.
  \item The EA maintains a \textbf{population} of candidate points. It wins
        when a sufficient fraction of its population (default 35\%)
        gathers inside the win radius.
  \item Four difficulty presets (\uielem{Easy}, \uielem{Normal},
        \uielem{Hard}, \uielem{Expert}) bundle an EA configuration with
        your probe \emph{reveal radius}: harder presets give the EA a
        bigger, more aggressive search and reveal less of the surface to
        you.
  \item After each of your probes you can open the \uielem{Replay} viewer
        to inspect exactly what the EA just did
        (Section~\ref{sec:replay}).
  \item If the EA solves the map first, you may study its solution and try
        to beat it in a second attempt.
\end{itemize}

\section{The EA settings panel}
\label{sec:ea-settings}

In Vs~EA mode you control the algorithm you are playing against. The
settings panel exposes the following parameters. The defaults listed below
correspond to the \uielem{Normal} difficulty preset, which is active when
the mode starts:

\begin{center}
\begin{tabularx}{\textwidth}{l l X}
  \toprule
  \textbf{Setting} & \textbf{Default} & \textbf{Meaning} \\
  \midrule
  Population size & 30 &
    Number of candidate points the EA maintains per generation. Larger
    populations explore more but converge more slowly. \\
  Max.\ generations & 200 &
    Upper limit on generations before the EA gives up. \\
  Crossover rate & 0.7 &
    Probability that two parents are recombined (instead of copied). \\
  Mutation rate & 0.3 &
    Probability that an offspring is mutated. \\
  Mutation strength & 0.25 &
    How far a mutation can move a point across the map. \\
  Mutation decay & 0.90 &
    Per-generation factor shrinking the mutation strength --- the EA takes
    big steps early and fine-tunes later. \\
  Win fraction & 0.35 &
    Fraction of the population that must sit inside the win radius for the
    EA to win. \\
  Selection & elitist &
    How parents are chosen: \emph{tournament}, \emph{roulette} or
    \emph{elitist}. \\
  Crossover & uniform &
    How two parents combine: \emph{arithmetic} (blend), \emph{uniform} or
    \emph{single-point}. \\
  Mutation & gaussian &
    Shape of the random perturbation: \emph{gaussian}, \emph{uniform} or
    \emph{cauchy} (occasional long jumps). \\
  \bottomrule
\end{tabularx}
\end{center}

Experimenting with these is the point of the game: try a huge mutation
strength and watch the population scatter; switch selection to
\emph{elitist} and watch it converge prematurely into a local minimum.

\section{The replay viewer}
\label{sec:replay}

The replay viewer is a full-screen overlay that replays the EA's last
generations \textbf{phase by phase}:

\begin{center}
  \emph{initial} $\to$ \emph{sorted} $\to$ \emph{elite} $\to$
  \emph{breeding} $\to$ \emph{mutating} $\to$ \emph{new generation} $\to$
  \emph{win check}
\end{center}

The left side shows the population as dots gliding across the map; the right
side lists every individual with its fitness and a role badge
(\uielem{ELITE}, \uielem{PARENT~A}, \uielem{PARENT~B}, \uielem{CHILD},
\uielem{WIN}). Use the playback controls to step forward/backward. A fitness chart tracks the population's best value across
generations.

\screenshot{PeakFinderEAReplay.png}{0.85}{The replay viewer (phase 1 of 7):
the population as dots on the map, ranked by fitness in the side list ---
every selection, crossover and mutation is made visible in the following
phases.}

\section{Sharing maps: map codes}
\label{sec:mapcodes}

Every map can be exported as a short text code (via the \uielem{Code}
dialog) and imported on any other device. Codes encode the minima positions,
the global minimum and the win radius. Saved maps and saved function
problems are also kept in a local sidebar for quick access.

\section{Benchmark functions}

Besides hand-built and random maps, the loader's \uielem{Functions} tab
offers a library of fourteen analytic benchmark functions, largely drawn
from the well-known BBOB test suite --- among them Sphere, Rosenbrock,
Ackley, Rastrigin, Schwefel, Eggholder and Gallagher's 101 Peaks. These
behave a little differently from maps: a function can have several winning
areas --- any spot that reaches the optimal value counts as a win, and those
spots can form larger regions or appear in more than one place. Each
function has its own code, so it can be loaded, saved and shared exactly
like a map.

\section{Function Tuner (developer tool)}

The \route{/FunctionTuner} page is a developer/instructor companion tool
for calibrating how the benchmark functions above are displayed: pick a
function, drag a slider to adjust its contrast (``sharpen'') exponent, and
watch the heatmap and win zone update live; a \uielem{Re-roll} button
sanity-checks the chosen value under the random rotations players see in
real play. Nothing on this page affects a running game --- it is not linked
from the user interface and only informs the shipped per-function display
settings.

% =====================================================================
\chapter{Maze Explorer}
% =====================================================================
\label{ch:maze}

\section{Concept}

Maze Explorer (\route{/MazeExplorer}) is a teaching exhibit built around a
simple contrast: first \emph{you} solve a maze the way an EA would --- by
writing a genome of moves and judging it with a fitness function --- then a
real EA solves the same kind of maze with a whole population at once, under
settings you control.

Mazes are either procedurally generated from a visible \textbf{seed} or
built by hand in Create mode. Because the seed pins down the exact maze, the
same maze can be re-run under different settings --- the key to the mode's
central experiment. The mode selector offers \uielem{Create},
\uielem{Solve} and \uielem{Control the EA}.

\section{Solve mode: you are the EA}

In \uielem{Solve} mode you play the role of a single individual. The maze is
hidden under fog-of-war, and you cannot walk it freely --- like an EA
genome, your solution is a \textbf{fixed-length string of moves} that you
write with the on-screen D-pad or the arrow keys
(\key{$\uparrow$}\,\key{$\downarrow$}\,\key{$\leftarrow$}\,\key{$\rightarrow$}).
Only when every slot is filled can you \uielem{Submit}: a walker then
follows your string through the fog, revealing the corridors it passes. A
move into a wall is simply \emph{wasted} --- the walker stands still for
that step.

Two \textbf{filmstrips} share one cursor: the top strip replays your last
submitted string (each move tinted by how close it landed to the goal), the
bottom strip holds your editable draft. After a run you edit individual
moves --- a hand-made \emph{mutation} --- and resubmit, iterating until the
walker reaches the goal. A selectable \textbf{fitness function}
(Section~\ref{sec:maze-fitness}) scores every run and repaints the strip
instantly, letting you see how the very same run is judged differently.

\screenshot{MazeSolveMode.png}{0.8}{Solve mode: a hand-written 63-move
string runs under fog-of-war. The filmstrips (top) hold the submitted and
the editing string, the D-pad (left) writes moves, and the fitness panel
(right) recolours the trail by closeness to the goal.}

\section{Control the EA (experiment mode)}

\uielem{Control the EA} is the heart of the game. Here an EA solves the maze
while you watch and dissect:

\begin{itemize}
  \item Each \textbf{individual} in the population is a complete
        \emph{path}: a fixed-length sequence of moves
        (\key{U}/\key{D}/\key{L}/\key{R}). A selectable \textbf{wall rule}
        decides what a blocked move does: \emph{waste} (stand still),
        \emph{break} (the walk ends at the crash), or \emph{repair} (the
        gene is fixed to a legal move, which is inherited).
  \item The maze sits behind a \uielem{Breed first generation} overlay
        while you tune the settings --- they shape generation~0. After
        that, each press of \uielem{Evolve} breeds the next generation;
        settings changed mid-run apply from the next generation and are
        pinned as markers on the fitness chart, so you can see exactly
        where a change bent the curve.
  \item The current \textbf{best individual} can be played back move by
        move on the maze, with the rest of the generation walking along as
        translucent \emph{ghosts} (toggleable): generation~1 is chaotic
        scribble; over generations the trails braid into a clean route to
        the goal.
  \item A per-generation \uielem{Dissect last generation} replay (like
        Peak Finder's, Section~\ref{sec:replay}) dissects selection,
        crossover and mutation. Crossover is especially vivid here: the
        child follows parent~A's route up to the splice cell, then
        continues on parent~B's.
\end{itemize}

\section{Fitness functions: the headline experiment}
\label{sec:maze-fitness}

The maze offers several \textbf{selectable fitness functions}, and comparing
them on the \emph{same seeded maze} is the game's main lesson:

\begin{description}[style=nextline]
  \item[Manhattan (naive objective)]
    Rewards ending close to the goal in straight-line (grid) distance,
    walls be damned. On most mazes this is \emph{deceptive}: it happily
    rewards a spot that is walled off from the goal, so the population
    piles up against the wall nearest the goal and gets stuck. Teaches:
    deceptive landscapes and premature convergence.
  \item[Geodesic (BFS field)]
    Rewards ending close to the goal \emph{through the corridors} (true
    walking distance). The same EA now climbs cleanly to the exit.
    Teaches: fitness shaping can fix a broken search.
  \item[Length penalty]
    Additionally rewards reaching the goal in fewer moves.
    Teaches: multi-objective tension between exploration and efficiency.
  \item[Novelty search]
    Ignores the goal entirely and rewards reaching cells that no previous
    individual has reached. Counter-intuitively, this often finds the exit
    when goal-chasing fails --- the classic Lehman \& Stanley maze result.
    Teaches: sometimes \emph{not} aiming at the objective works better.
\end{description}

\textbf{Suggested experiment:} note the maze seed, run the EA with Manhattan
fitness and watch it get trapped; then re-run the identical maze with
geodesic fitness, then with novelty search. Same maze, same algorithm ---
three completely different behaviours. In Solve mode the same selector
exists, so you can also feel the difference on your \emph{own} runs ---
switching repaints the filmstrip without resubmitting.

\section{Create mode}

\uielem{Create} mode is a maze editor: draw and erase walls on the grid,
place the start and the goal, and save the result under a name. Saved mazes
are kept locally and can be shared via codes; from the editor (or the setup
screen) a maze is handed straight to Solve or Control-the-EA mode, and an
\uielem{Edit this maze} button in experiment mode brings the current maze
back into the editor.

% =====================================================================
\chapter{Shooter Game}
% =====================================================================
\label{ch:shooter}

\section{Concept}

The Shooter Game (\route{/ShooterGame}) is a top-down arena shooter with a
twist: the enemies are controlled by a \textbf{genetic algorithm}. Each
enemy agent's behaviour is encoded as a strand of \textbf{DNA} --- eight
genes controlling aggression, dodging, accuracy, preferred range, speed,
aim prediction, fire rate and bullet speed. Between rounds the population
\emph{evolves}, using a recording (``ghost'') of your previous round as
training data. The enemies do not just get harder --- they adapt to
\emph{how you play}.

\section{Controls}

\begin{center}
\begin{tabular}{ll}
  \toprule
  \textbf{Action} & \textbf{Input} \\
  \midrule
  Move            & \key{W}\,\key{A}\,\key{S}\,\key{D} or arrow keys
                    (desktop) / virtual joystick (touch) \\
  Aim             & mouse (desktop) / aim zone (touch) \\
  Shoot           & left mouse button or \key{Space} (desktop) /
                    aim zone (touch) \\
  \bottomrule
\end{tabular}
\end{center}

There are no further bindings --- no dash or similar special moves.

\section{Game modes}

The shooter lobby (\route{/lobby/shooter}) leads to three modes, each with
its own \uielem{?}~help dialog in-game:

\subsection{Solo}

Round-based duels against evolved agents: land more hits than you take
within the 20-second round to win it. After every round the population
evolves against your ghost recording; the on-screen \uielem{DNA display}
shows the current opponent's genes, so you can literally watch traits like
\emph{aggression} or \emph{accuracy} rise in response to your tactics.
Difficulty presets decide the opponents' starting DNA and how many
generations they secretly pre-train before round one. The
\route{/Analytics} page keeps per-round records of your Solo runs.

\subsection{Horde}

A survival mode on obstacle maps --- and a showcase of a
\textbf{steady-state} EA: there are no rounds or waves. Every time you kill
an agent, its replacement is bred on the spot from the survivors, so the
swarm adapts \emph{continuously} for as long as you stay alive. Agents here
carry extra genes (a cyclic movement pattern, body size and opacity ---
smaller, fainter agents are genuinely harder to hit, creating real
selection pressure). Obstacles matter: solid-bordered ones block bullets
--- duck behind them for cover --- while dashed ones only block movement.
During a run you collect stackable \textbf{mods} (power-ups) that modify
your stats or your shot pattern. Maps are selectable, and a companion
\textbf{map editor} (\route{/HordeMapEditor}) lets you build and play your
own.

\subsection{Community Raidboss}

A shared, community-evolved opponent: one population of boss DNA lives in
the cloud (Firebase). Starting a raidboss session claims one individual
from that population; your round against it scores its fitness, and the
result flows back into the shared pool, which evolves from everyone's
fights. Every player therefore trains the same adversary --- the lobby
shows how many contributors have shaped it so far.

\section{Reading the DNA display}

Each agent's genome is a vector of eight genes, all in $[0,1]$:

\begin{center}
\begin{tabularx}{\textwidth}{l X}
  \toprule
  \textbf{Gene} & \textbf{Effect} \\
  \midrule
  Aggression       & How strongly the agent chases you. \\
  Dodge weight     & How hard it tries to avoid your bullets. \\
  Shoot accuracy   & Aim precision (1 = perfect aim). \\
  Preferred range  & The combat distance it tries to maintain. \\
  Movement speed   & Speed bonus on top of the base speed. \\
  Predict lead     & How far it leads its shots ahead of your movement. \\
  Fire rate        & How quickly it shoots. \\
  Bullet speed     & How fast its bullets travel. \\
  \bottomrule
\end{tabularx}
\end{center}

(Horde agents extend this genome with a cyclic movement pattern, body size
and opacity --- see the Horde mode above.)

Watching these values drift over rounds is the game's evolutionary
storytelling: camp in a corner and \emph{accuracy} and \emph{range} genes
win; rush constantly and \emph{dodge weight} rises.

\section{How the evolution works (in brief)}

After each round every agent receives a fitness score based on its
performance against you (damage dealt, survival, \ldots). The algorithm
then applies tournament selection, single-point crossover and per-gene
mutation to produce the next generation. Before your first round, the
population can be pre-trained by simulating agents against each other, so
even round~1 opponents are not purely random.

% =====================================================================
\chapter{A Small EA Glossary}
% =====================================================================
\label{ch:glossary}

Plain-language definitions of the terms used across all three games.

\begin{description}[style=nextline]
  \item[Individual] One candidate solution. A point on the map (Peak
    Finder), a move sequence (Maze), or a DNA vector (Shooter).
  \item[Population] The set of individuals alive in one generation.
  \item[Fitness] A number scoring how good an individual is. In Peak Finder
    lower is better (distance-like); in the Shooter higher is better.
  \item[Generation] One full cycle of evaluate $\to$ select $\to$ breed
    $\to$ mutate.
  \item[Selection] Choosing which individuals get to reproduce.
    \emph{Tournament}: pick a few at random, the best of them wins.
    \emph{Roulette}: chance proportional to fitness.
    \emph{Elitist}: only the best reproduce.
  \item[Crossover] Combining two parents into a child.
    \emph{Arithmetic}: blend their values. \emph{Single-point}: take the
    first part from one parent, the rest from the other. \emph{Uniform}:
    coin-flip per gene.
  \item[Mutation] Randomly perturbing a child. Keeps the population
    exploring instead of collapsing onto one solution.
  \item[Elitism] Copying the best few individuals unchanged into the next
    generation so progress is never lost.
  \item[Local optimum] A solution better than all its neighbours but worse
    than the global best --- where greedy search gets stuck.
  \item[Deceptive landscape] A problem where following the apparent
    gradient leads \emph{away} from the true optimum (see the Euclidean
    maze fitness, Section~\ref{sec:maze-fitness}).
  \item[Premature convergence] The whole population collapsing onto one
    (often local) optimum before exploring enough.
  \item[Novelty search] Selecting for \emph{new behaviour} instead of
    fitness --- rewarding individuals that reach places no one has reached
    before.
\end{description}

% =====================================================================
\chapter{Troubleshooting \& FAQ}
% =====================================================================

\begin{description}[style=nextline]
  \item[The page loads but a game screen stays blank.]
    Make sure your browser is up to date (the app uses modern Web Worker
    and Canvas APIs). Try a hard reload
    (\key{Ctrl}+\key{Shift}+\key{R} / \key{Cmd}+\key{Shift}+\key{R}).
  \item[The EA in Vs~EA mode never seems to move.]
    The EA only advances when \emph{you} probe --- it runs a fixed number
    of generations per probe. Place a probe, then open the replay viewer.
  \item[Hints stopped appearing.]
    First check that the light-bulb toggle is on --- it is remembered
    between visits. Peak Finder and Maze Explorer hints re-appear whenever
    their trigger occurs; only a few one-time hints (e.g.\ the shooter
    lobby tour invitations) are shown once per session --- the
    \uielem{Reset} button next to the hint toggle re-arms those, and the
    lobby tours can always be re-run via \uielem{Take the Tour}.
  \item[My saved maps / mazes disappeared.]
    Saved content is stored in your browser's local storage. It does not
    transfer between browsers or devices, and clearing site data deletes
    it. Use map codes (Section~\ref{sec:mapcodes}) to back up or transfer
    maps.
  \item[Performance is poor on my machine.]
    Reduce the EA population size in the settings panel, and in the maze
    experiment turn off the ghost walkers. The EAs themselves run in
    background Web Workers, so the UI should stay responsive even with
    heavy settings.
  \item[Does the app collect data?]
    The app runs client-side and requires no account. Only the community
    Raidboss mode communicates with a server (Firebase): it stores the
    shared boss population --- DNA vectors with their fitness statistics
    and a contributor counter. No personal data is transmitted.
\end{description}

\appendix

% =====================================================================
\chapter{Route Reference}
% =====================================================================

\begin{center}
\begin{tabularx}{\textwidth}{l X}
  \toprule
  \textbf{Route} & \textbf{Page} \\
  \midrule
  \route{/}                & Home / welcome page \\
  \route{/Dashboard}       & Game selection dashboard \\
  \route{/PeakFinder}      & Peak Finder (Chapter~\ref{ch:peakfinder}) \\
  \route{/MazeExplorer}    & Maze Explorer (Chapter~\ref{ch:maze}) \\
  \route{/ShooterGame}     & Shooter Game (Chapter~\ref{ch:shooter}) \\
  \route{/lobby/shooter}   & Shooter lobby / mode selection \\
  \route{/HordeGame}       & Horde mode \\
  \route{/HordeMapEditor}  & Horde map editor \\
  \route{/FunctionTuner}   & Function display tuning tool (developer tool,
                             not linked from the UI) \\
  \route{/Analytics}       & Per-round analytics of Solo shooter runs \\
  \bottomrule
\end{tabularx}
\end{center}

\end{document}
